\chapter{TESTING AND EVALUATION}
[Placeholder for Capstone 1: this chapter will be completed after the system is built and tested.]

\section{Testing and Evaluation}
This chapter presents the testing and evaluation of \SystemName. It outlines the procedures used to test the functionality, performance and usability of the system.

The study tested \SystemName\ with ten (10) respondents: five (5) I.T. professionals and five (5) [staff of office].

\section{Test Case Result}
\begin{table}[H]\caption{Module 1 - Login}\centering
\begin{tabularx}{\textwidth}{|>{\raggedright\arraybackslash}p{1.1in}|Y|>{\centering\arraybackslash}p{0.6in}|>{\centering\arraybackslash}p{0.6in}|}\hline
\hd{STEPS} & \hd{TEST STEPS} & \hd{PASS} & \hd{FAIL}\\ \hline
Login & To log in to the system, input username and password and click the login button. & [10] & [0]\\ \hline
[Function 2] & [Steps the tester follows.] & [10] & [0]\\ \hline
\end{tabularx}\end{table}

The table above shows the result of the testing from ten (10) respondents that mark the functionality of Module 1 - Login. It features [what the module does].

[Repeat for Module 2, Module 3 and so on until every module is covered.]

\section{Analysis and Result}
During testing, every system function relevant to the target users was included. This directly benefits \OfficeName\ by addressing the inefficiencies in its manual processes.

The feedback from the ten (10) respondents included [recommendation]. [Issues found and whether they affected use.] In summary, the test case assessments showed that [what the system does well and what can be improved], and the users found that the features meet their basic needs.
